\section{Open Questions and Recommended Next Steps}

The five open questions below are genuine reproducibility/extension gaps left by this replication. Each is paired with a concrete, free-endpoint-only probe.

\subsection*{Q1 --- Independent MC re-derivation of the fitted coefficients}
\textbf{Question.} Can the Kundr\'at 2020 fitted coefficients (Tables~1--2) be independently re-derived from a public track-structure code such as Geant4-DNA (Opt2/Opt4 physics list) or TOPAS-nBio, within the paper's quoted $<2\%$ RMS envelope (up to $\sim 9\%$ for DSB clusters)?

\textbf{Basis.} PARTRAC is closed in-house GSF/HMGU code and the raw simulation-point tables behind Figs.~1--5 are not deposited anywhere accessible. Without an independent MC channel, the paper's central validation claim is not falsifiable outside GSF.

\textbf{Next steps.} Set up Geant4-DNA Option~2 with an equivalent chromatin-fibre geometry; run H, He, C, Ne on the same LET grid; refit Eqs.~1--2 and compare coefficient-by-coefficient. Free, feasible on m1 or a single Spark GPU node in a few CPU-days.

\subsection*{Q2 --- Extrapolation to non-standard track structures}
\textbf{Question.} How do the analytical formulas extrapolate to heavy ions past Ne (Ar, Fe at deep-space or hadron-therapy energies) and to low-energy secondary electrons (sub-keV)?

\textbf{Basis.} The paper warns against out-of-range use; our own analysis surfaced a concrete failure --- the H-proton DSB-site row lacks the overkill denominator and diverges past the fitted range, contradicting the $\sim$15 sites/Gy/Gbp peak.

\textbf{Next steps.} Encode the per-ion fitted LET domain as hard clamps in \texttt{formulas.py}; generate Geant4-DNA points for Ar and Fe at 1--500~keV/\textmu m and see whether Eq.~1/Eq.~2 with refitted coefficients still describe the curves; probe sub-keV electrons with TOPAS-nBio's low-energy EM list; publish a fit-domain map.

\subsection*{Q3 --- Integration with modern chromatin geometry models}
\textbf{Question.} Can these formulas be integrated with modern chromatin-geometry models (fractal globule, Hi-C-derived nuclear geometry, phase-separated compartments) to produce spatially-resolved DSB-cluster predictions, or are the fits tied to the PARTRAC-era 30~nm fibre representation?

\textbf{Basis.} DSB-site classification depends on the chromatin model inside PARTRAC. Modern data-driven geometries produce different local-density distributions, which change DSB-clustering statistics for identical track structures.

\textbf{Next steps.} Replace the implicit uniform-fibre geometry with a voxelized chromatin density map derived from Micro-C or 3D-genome models; recompute DSB-site clustering (DBSCAN, 25 bp threshold) on the same track set; report DSB-site sensitivity to chromatin choice.

\subsection*{Q4 --- Applicability to FLASH dose-rate regimes}
\textbf{Question.} Do the fitted formulas remain quantitatively accurate at ultra-high dose rates (FLASH, $\geq 40$~Gy/s), where inter-track chemistry and radical recombination matter, or do they systematically over-predict DSB yields?

\textbf{Basis.} Kundr\'at's PARTRAC runs were in the single-track / low-dose-rate regime. FLASH operates in a regime where oxygen depletion, radical recombination and inter-track effects are hypothesized to reduce net damage --- none of which is in Eqs.~1--2.

\textbf{Next steps.} Compare Eq.~1/Eq.~2 predictions against published FLASH-regime yield measurements (Adrian~2020, Wilson~2020) and against dose-rate-aware MC codes (TRAX-CHEM, GEANT4-DNA-Chem) at 0.03~Gy/s vs 100~Gy/s; report the $Y_{\text{FLASH}}/Y_{\text{conv}}$ ratio.

\subsection*{Q5 --- Sensitivity to functional form; ML-driven formula discovery}
\textbf{Question.} How sensitive is the $<2\%$ RMS headline to the specific functional forms in Eqs.~1--2, and could ML-driven symbolic regression (PySR, AI-Feynman) discover more parsimonious or physically-motivated surrogates for other track-structure MC codes (Geant4-DNA, TOPAS-nBio, KURBUC)?

\textbf{Basis.} Eqs.~1--2 are hand-crafted phenomenological forms; there is no a-priori argument that log-squared damping (Eq.~1) or power-law-over-power-law (Eq.~2) is the uniquely correct compression. Different functional forms could fit equally well and extrapolate more safely.

\textbf{Next steps.} Generate a yield grid with an open MC; run PySR with a physics-friendly operator set; report top-$k$ expressions with RMS and parameter count; score in-range fit and out-of-range behaviour; assess interpretability. Free, $\sim$1 GPU-day per ion on a Spark node.
